\documentclass[aspectratio=169,10pt]{beamer}
\newcommand{\reporoot}{../..}
% Shared preamble of the main deck and the topic decks.
% Define \reporoot (the path from the .tex file to the repository root)
% before \input-ing this file.
\usetheme[progressbar=frametitle,numbering=fraction,sectionpage=progressbar]{metropolis}

\usepackage{amsmath,amssymb}
\usepackage{booktabs}
\usepackage{array}
\newcolumntype{L}[1]{>{\raggedright\arraybackslash}p{#1}}
\usepackage{graphicx}
\usepackage{tikz}
\usetikzlibrary{arrows.meta,positioning,decorations.pathreplacing}
\usepackage[authoryear,round]{natbib}
\usepackage{xcolor}

\graphicspath{{\reporoot/figures/}}

% Every simulated number on the slides is produced by experiments/*.py and read from here.
\input{\reporoot/results/e01_estimators.tex}
\input{\reporoot/results/e02_scaling.tex}
\input{\reporoot/results/e03_var_backtests.tex}
\input{\reporoot/results/e04_es_backtest.tex}
\input{\reporoot/results/e05_pla.tex}
\input{\reporoot/results/e06_estimation_risk.tex}
\input{\reporoot/results/e07_model_uncertainty.tex}
\input{\reporoot/results/e08_asrf.tex}
\input{\reporoot/results/e09_copula_is.tex}
\input{\reporoot/results/e10_calibration_tests.tex}
\input{\reporoot/results/e11_discrimination.tex}
\input{\reporoot/results/e12_exposure.tex}
\input{\reporoot/results/e13_cva_wwr.tex}
\input{\reporoot/results/e14_nested.tex}
\input{\reporoot/results/e15_variance_reduction.tex}
\input{\reporoot/results/e16_lda.tex}
\input{\reporoot/results/e17_evt.tex}
\input{\reporoot/results/e18_reverse_stress.tex}
\input{\reporoot/results/e19_regimes.tex}
\input{\reporoot/results/gallery.tex}

\newcommand{\E}{\mathbb{E}}
\newcommand{\Var}{\operatorname{Var}}
\newcommand{\VaR}{\mathrm{VaR}}
\newcommand{\ES}{\mathrm{ES}}
\newcommand{\dd}{\mathrm{d}}
\definecolor{accent}{HTML}{0C5DA5}
\definecolor{warn}{HTML}{C0392B}
\newcommand{\hl}[1]{\textcolor{accent}{\textbf{#1}}}

\setbeamerfont{footnote}{size=\tiny}
\setbeamerfont{caption}{size=\scriptsize}

\author{Yaohui Shu}
\institute{Stony Brook University}
\date{November 2026}

% Topic decks: a references frame that lists only the works cited in the deck.
\newcommand{\deckreferences}{%
  \begin{frame}[allowframebreaks]{References}
    \tiny
    \renewcommand{\bibsection}{}
    \bibliographystyle{plainnat}
    \bibliography{\reporoot/slides/refs}
  \end{frame}}

% Topic decks: \decknumber{n} puts the deck number on the title page.
\newcommand{\decknumber}[1]{\institute{Stony Brook University\hfill{\scriptsize\color{gray}Deep dive #1 of 9}}}


\title{Model Risk and Estimation Risk}
\subtitle{SR 26-2, parameter uncertainty, and how wrong a correct model can be}
\decknumber{3}

\begin{document}

\maketitle

\begin{frame}{The question}
  \small
  \begin{columns}[T]
    \begin{column}{0.55\textwidth}
      Every risk number comes from a model, and every model is fitted to limited data. \hl{How much of a risk estimate is the risk, and how much is the model and the sample?}
      \begin{itemize}
        \item Estimation risk: the right model with estimated parameters.
        \item Model uncertainty: which model is right at all.
      \end{itemize}
      Both can be quantified by simulation and closed forms, and both are what validation teams are asked to challenge.
    \end{column}
    \begin{column}{0.42\textwidth}
      \begin{alertblock}{In this deck}
        \small
        \begin{enumerate}
          \item The 2026 US guidance, SR 26-2.
          \item Exceedance probability of estimated VaR.
          \item Sampling error of ES.
          \item Worst cases and challenger models.
        \end{enumerate}
      \end{alertblock}
    \end{column}
  \end{columns}
\end{frame}

\begin{frame}{Estimation error and worst cases, seen directly}
  \centering
  \includegraphics[width=0.95\textwidth,height=0.54\textheight,keepaspectratio]{gallery/g03_model_risk.png}\par\vspace{0.4em}
  \footnotesize\raggedright
  (a) Sampling distribution of the plug-in normal 99\% VaR over the true VaR for window lengths $n$ (400,000 samples each), with the exact probability that the next loss exceeds it. (b) Worst-case loss distributions within relative entropy $\eta$ of a standard normal: each is the normal shifted by $\sqrt{2\eta}$, and its 97.5\% ES (dots) rises by the same amount.
\end{frame}

\begin{frame}{SR 26-2: revised guidance on model risk management}
  \small
  Issued 17 April 2026 by the Federal Reserve, OCC and FDIC; it \hl{supersedes SR 11-7} (2011) and SR 21-8 \citep{fed2026sr262}.
  \begin{columns}[T]
    \begin{column}{0.5\textwidth}
      \begin{itemize}
        \item Most relevant to banking organizations with over \$30 billion in total assets; risk-based and tailored.
        \item A model is a ``complex quantitative method, system, or approach'' that uses statistical, economic or financial theory to turn input data into quantitative estimates. Simple arithmetic and deterministic rules are out.
        \item Generative and agentic AI models are outside its scope; other AI models are in.
      \end{itemize}
    \end{column}
    \begin{column}{0.46\textwidth}
      \begin{itemize}
        \item Model risk depends on inherent risk, exposure, purpose and use; exposure and purpose determine materiality.
        \item Validation components: conceptual soundness, outcomes analysis (including back-testing), ongoing monitoring.
        \item Effective challenge: critical analysis by objective experts with the standing to effect change.
      \end{itemize}
    \end{column}
  \end{columns}
  \vspace{0.3em}
  This deck makes two of these concrete: outcomes analysis has sampling error, and conceptual soundness needs alternatives to compare against.
\end{frame}

\begin{frame}{Estimated VaR is exceeded more often than its level says}
  \small
  \begin{columns}[T]
    \begin{column}{0.52\textwidth}
      \textbf{Normal losses, plug-in VaR.} From $n$ i.i.d.\ normal losses with sample mean $\bar L$ and standard deviation $s$, the next loss satisfies
      \[ \frac{L_{n+1} - \bar L}{s\sqrt{1 + 1/n}} \sim t_{n-1}, \]
      so the plug-in VaR $\bar L + s\,z_\alpha$ is exceeded with probability
      \[ P = 1 - T_{n-1}\!\left(\frac{z_\alpha}{\sqrt{1+1/n}}\right) > 1 - \alpha. \]
      The predictive VaR $\bar L + s\sqrt{1+1/n}\;t^{-1}_{n-1}(\alpha)$ has exact coverage.
    \end{column}
    \begin{column}{0.45\textwidth}
      \textbf{Historical simulation.} If VaR is the $k$-th smallest of $n$ losses, then for any continuous distribution
      \[ \E[P(L_{n+1} > L_{(k)})] = \frac{n + 1 - k}{n + 1}, \]
      because $F(L_{(k)})$ is the $k$-th uniform order statistic. With $n = 250$ and $k = 248$: $\ESixHsTwoFifty\%$, not 1\%.
      \vspace{0.3em}

      Both are standard results; this repository checks them by simulation.
    \end{column}
  \end{columns}
\end{frame}

\begin{frame}{Estimation risk in numbers}
  \centering
  \includegraphics[width=0.86\textwidth]{fig06_estimation_risk.png}
  \begin{columns}[T]
    \begin{column}{0.5\textwidth}
      \footnotesize
      (a) Plug-in normal 99\% VaR: exceeded $\ESixPlugTwentyFive\%$ of the time with $n = 25$ and $\ESixPlugTwoFifty\%$ with $n = 250$ (Monte Carlo $\ESixMcTwoFifty\%$); predictive VaR $\ESixPredTwoFifty\%$. HS with the $\lceil 0.99n\rceil$-th order statistic has a sawtooth: $\ESixHsTwoHundred\%$ at $n = 200$.
    \end{column}
    \begin{column}{0.46\textwidth}
      \footnotesize
      (b) Empirical ES$_{97.5}$ from 250 days: relative standard error $\ESixSdNorm\%$ for normal losses and $\ESixSdT\%$ for $t_3$ (bias $\ESixBiasT\%$). 95\% bootstrap intervals cover the truth $\ESixCoverNorm\%$ and $\ESixCoverT\%$ of the time.
    \end{column}
  \end{columns}
\end{frame}

\begin{frame}{Model uncertainty: worst cases and challengers}
  \small
  \begin{columns}[T]
    \begin{column}{0.52\textwidth}
      \textbf{Relative-entropy worst case} \citep{glasserman2014robust}. Among models $Q$ with $D(Q\,\|\,P) = \E_Q[\ln(\dd Q/\dd P)] \le \eta$, the largest $\E_Q[V]$ is attained by the exponential tilt
      \[ \frac{\dd Q^*}{\dd P} = \frac{e^{\theta V}}{\E_P[e^{\theta V}]}, \]
      with $\theta \ge 0$ chosen so that the divergence equals $\eta$. It needs only simulated scenarios from the nominal model: reweight them.
      \vspace{0.2em}

      Here $V = L$: the model that maximizes expected loss, then its ES. For $L \sim \mathcal N(0,1)$ the tilt is $\mathcal N(\theta, 1)$ with $\theta = \sqrt{2\eta}$, so ES rises by exactly $\theta$. For heavy tails $\E[e^{\theta L}] = \infty$: the tilt does not exist.
    \end{column}
    \begin{column}{0.45\textwidth}
      \textbf{Challenger models.} Validation compares the champion with alternatives fitted to the same data. The spread across reasonable models is a direct measure of model risk for that number.
      \vspace{0.3em}

      Here: the five estimators of deck 1, each applied to the same 500-day GARCH-t window, for 1,000 windows.
    \end{column}
  \end{columns}
\end{frame}

\begin{frame}{How large is model risk?}
  \centering
  \includegraphics[width=0.94\textwidth]{fig07_model_uncertainty.png}
  \footnotesize
  \begin{columns}[T]
    \begin{column}{0.32\textwidth}
      (a) Normal losses, 2 million reweighted scenarios: within $\eta = 0.1$ the worst-case ES is $\ESevenWcNormTenth$ times nominal; the largest gap to the exact $\ES + \sqrt{2\eta}$ is $\ESevenWcErr$.
    \end{column}
    \begin{column}{0.32\textwidth}
      (b) Student $t_4$ losses, $\eta = 0.05$: the ``worst case'' is $\ESevenTSmallLo$ to $\ESevenTSmallHi$ times nominal with $10^4$ scenarios and $\ESevenTLargeLo$ to $\ESevenTLargeHi$ with $4\times10^6$. It depends on the largest draws, not on the model.
    \end{column}
    \begin{column}{0.32\textwidth}
      (c) Largest over smallest ES$_{97.5}$ among five estimators on the same window: median $\ESevenSpreadMed$, 90th percentile $\ESevenSpreadNinety$; above 1.5 in $\ESevenSpreadOver\%$ of windows.
    \end{column}
  \end{columns}
\end{frame}

\begin{frame}{Discussion and takeaways}
  \small
  \begin{columns}[T]
    \begin{column}{0.5\textwidth}
      \textbf{What it means for validation}
      \begin{itemize}
        \item Outcomes analysis has its own sampling error: a correctly specified model fitted to 250 days misses its nominal level, and ES estimates from one year carry tens of percent of error.
        \item Bootstrap intervals for ES from one year cover the truth less often than stated, even for normal losses.
        \item Model uncertainty can be bounded without a second model (relative entropy), but only for light tails; for heavy tails, measure it against challengers.
      \end{itemize}
    \end{column}
    \begin{column}{0.46\textwidth}
      \textbf{For a model risk team}
      \begin{enumerate}
        \item Quote risk numbers with an estimation error, not as points.
        \item Use predictive (not plug-in) quantiles when the sample is short.
        \item Keep at least one challenger per material model and track the spread.
        \item Express judgmental overlays as a divergence budget, where the tails allow it.
      \end{enumerate}
      \textbf{Open questions.} How to set $\eta$ from data and governance, not by hand; worst-case bounds that stay finite for heavy-tailed losses.
    \end{column}
  \end{columns}
\end{frame}

\begin{frame}[fragile]{Reproduce}
  \small
  Code: \texttt{model\_risk.py}. Experiments: \texttt{e06\_estimation\_risk.py}, \texttt{e07\_model\_uncertainty.py}.
\begin{verbatim}
from stochastic_risk import model_risk as mr
mr.plugin_normal_exceedance(n=250, alpha=0.99)    # exact exceedance probability
mr.predictive_normal_var(losses, 0.99)            # exact coverage
mr.hs_expected_exceedance(n=250, k=248)           # distribution-free
mr.worst_case_es(losses, alpha=0.975, eta=0.1)    # Glasserman-Xu tilt (light tails)
\end{verbatim}
\end{frame}

\deckreferences
\end{document}
